\chapter{Building project 1: oracle vs zero-leakage IBM}\label{ch:build1}

\section{Layout}
\begin{lstlisting}[style=sh]
sexconflict/
  config.py                 # conditions, seeds (MASTER_SEED + seed_for(name)), threads
  deterministic/            # ORACLES
    poster_model.py         #   3-D map, sympy Jacobian, closed-form B, fate(), fate_fast()
    owen_parsons.py         #   recursions, latent roots, equilibria
    mothers_curse.py        #   cytonuclear recursion
    sa_classics.py          #   Kidwell/Rice/Fry, Kimura, full diffusion, Kolmogorov, Rice 1987
  ibm/                      # SIMULATORS (zero leakage)
    kernels.py              #   numba: haploid 2-locus, diploid+mito, X-linked, XY
    opencl_backend.py       #   same primitives on GPU
    runners.py              #   founders, burn-in, measurement
  analysis/stats.py, ml.py  # Wilson CI, ratio tests, KS, logistic fits; GP/MLP/SBI
  experiments/*.py          # one module per suite -> figures, Parquet
scripts/run_all.py, make_report.py
tests/                      # 36 tests incl. zero-leakage
\end{lstlisting}

\section{Step 1: the poster oracle}\label{sec:poster}
\begin{stepbox}[Reconstruct an under-specified model]
The poster gives the fitness table and $\pstar$ but not the life cycle. Write down candidate life cycles, derive $\pstar$ for each symbolically, and keep the one that reproduces the published formula exactly. Here that is: haploid selection in each sex, then random union, then meiosis with recombination $r$.
\end{stepbox}

Haplotypes are $x=(A_1M_1,A_1M_2,A_2M_1,A_2M_2)$. Fitnesses: males $(1+a,1+a,1,1)$; females $(1+bk_1,1,1+bk_2,1+b)$.
\begin{equation}
x'_{ij}=\tfrac12(1-r)(f_{ij}+m_{ij})+\tfrac12 r\,(f_{A_i}m_{M_j}+m_{A_i}f_{M_j}),
\end{equation}
where $m_{ij}=x_{ij}w^m_{ij}/\bar w_m$, $f_{ij}=x_{ij}w^f_{ij}/\bar w_f$, and $f_{A_i}$, $m_{M_j}$ are allele marginals after selection.

\textbf{Implementation steps}
\begin{enumerate}
\item Write the map vectorised over a leading batch axis (numpy broadcasting), so whole $(a,r)$ grids iterate at once.
\item Derive the Jacobian with sympy. Check it against finite differences ($<10^{-10}$) and against the closed form ($10^{-12}$).
\item \textbf{Invasion.} The face $x_{11}=x_{21}=0$ is invariant, so the Jacobian at $(0,\pstar,0)$ is block-triangular. M$_1$ invades iff $\rho(B)>1$, with $B$ the $2\times2$ block given in the expert report.
\item \textbf{Global fate} (\code{fate}): iterate until convergence with per-element stopping. A numba version (\code{fate_fast}) is about 8$\times$ faster; regression-test it against a converged numpy reference.
\item \textbf{Thresholds} by bisection on the initial M$_1$ dose; \textbf{basins} from Dirichlet(1,1,1,1) initial conditions.
\end{enumerate}

\section{Step 2: the IBM kernels}
\textbf{One generation} of every kernel is built from events on individuals:
\begin{enumerate}
\item Assign sex: a fair coin, or the inherited Y in the XY kernel.
\item Survival: individual $i$ survives with probability $w_{\text{sex},g_i}/\max_g w_{\text{sex},g}$.
\item Each of $N$ offspring picks a uniformly random surviving mother and father.
\item Transmission: Mendelian segregation, crossover with probability $r$, maternal cytoplasm, Y to sons.
\end{enumerate}
The haploid two-locus kernel, slightly shortened:
\begin{lstlisting}
@njit(cache=True)
def _haploid_one(A, M, n_gen, wtab, r, stop_on_m_absorb, rec):
    N = A.shape[0]
    ...
    for g in range(n_gen + 1):
        for k in range(4):                       # measurement only: count
            rec[g, k] = 0
        for i in range(N):
            rec[g, 2 * A[i] + M[i]] += 1
        if g == n_gen:
            break
        _assign_sex(sex)                         # biology
        for i in range(N):
            cls[i] = 2 * A[i] + M[i]
        nf, nm = _survivors(sex, cls, wtab, fem, mal)
        for j in range(N):
            mom = fem[np.random.randint(nf)]
            dad = mal[np.random.randint(nm)]
            if np.random.random() < 0.5:         # which parent gives locus A
                p, q = mom, dad
            else:
                p, q = dad, mom
            A2[j] = A[p]
            M2[j] = M[q] if np.random.random() < r else M[p]   # crossover
        A[:] = A2; M[:] = M2

@njit(parallel=True, cache=True)
def haploid_ensemble(init_counts, N, n_gen, wtab, r, n_rep, seed, stop_on_m_absorb):
    out = np.zeros((n_rep, n_gen + 1, 4), np.int64)
    for rep in prange(n_rep):
        np.random.seed(seed + rep)               # per-replicate reseed (numba RNG)
        ...
\end{lstlisting}
The other kernels are \code{diploid_ensemble} (with maternal mito), \code{xlinked_ensemble} and \code{xy_ensemble} (genetic sex determination: male iff the Y was inherited from the father).

\section{Step 3: confrontation protocols}
\begin{tabularx}{\textwidth}{lXX}\toprule
ID & Oracle & \IBM\ and comparison\\\midrule
P0 & $\pstar$ & burn-in from 50\%; time average; RMSE vs $N$\\
P2 & $\rho(B)$ and trajectory from $(\pstar, 2\%)$ & burn-in, then mutate 2\% to M$_1$, then 300 generations; sign of mean change (95\% CI) vs oracle\\
P4 & bisection threshold & $P(\text{resolution})$ vs dose at 3 $N$; logistic midpoint and width\\
P5 & basin fractions & the same initial conditions as founders; per-condition agreement\\
P6 & trajectory & between-population SD $\propto N^{-1/2}$\\
O1--O4 & latent roots, equilibria, flow field & establishment maps, equilibria, basin agreement\\
C1--C4 & cytonuclear recursion & KS on endpoints, thresholds, $P_\mathrm{fix}$, arms race\\
E1--E5 & Kidwell/Fry/Rice, Kimura, Rice 1987, Kolmogorov & rare-invasion maps, fixation, XY kernel, persistence\\\bottomrule
\end{tabularx}

\noindent\textbf{Statistics helpers} (\code{analysis/stats.py}):
\begin{itemize}
\item Wilson intervals for proportions;
\item a ratio test of early spread ($\log_2$ of mean frequency change);
\item KS tests;
\item logistic fits of $P(\text{event})$ against dose (midpoint and width);
\item RMSE.
\end{itemize}

\section{Step 4: mother's curse model}\label{sec:curse}
The state is the zygote distribution $Z[c,g]$: mito type $c$ (maternal) and restorer copies $g$. Viabilities are
\[
\text{females: }(1+s_fc)(1-\kappa d_g),\qquad \text{males: }(1-s_mc(1-\rho d_g))(1-\kappa d_g),\qquad d=(0,h,1).
\]
The \IBM\ kernel carries a maternal mito flag per individual. The tests are the strong form (growth $1+s_f$ independent of $s_m$), the twofold rule, the weak form ($P_\mathrm{fix}$ equals the initial frequency) and the restorer arms race.

\section{Step 5: extended suite}
\begin{itemize}
\item \textbf{Protected-polymorphism maps.} Two rare-invasion experiments per $(s_f,s_m)$ cell, autosomal and X; polymorphism is protected iff both rare alleles invade.
\item \textbf{Kimura vs full diffusion.} Compute $M(p)$ from the full sex-specific recursion and integrate $\psi$ numerically. The haploid form with $N_e=N/2$ applies to mito.
\item \textbf{Rice 1987.} A new recursion tracks the SA allele on X in eggs, on X-sperm and on Y-sperm, with crossover $r$ to the sex-determining region; it reduces to the autosomal model at $r=\tfrac12$ (tested).
\item \textbf{Persistence times.} Finite differences for $\tfrac12VT''+MT'=-1$ on 4,001 nodes.
  \begin{itemize}
  \item \textbf{$N_e$ tuner:} measure $N_e$ from the \IBM's own one-step variance, $\mathrm{Var}(\Delta p)=pq/(2N_e)$.
  \item Random survival gives $N_e/N\approx0.9$, which cut the error from 11.5\% to 3.7\%.
  \end{itemize}
\end{itemize}

\section{Step 6: GPU backend (OpenCL, optional)}
\begin{itemize}
\item On Windows with AMD Polaris (RX 560X), ROCm/HIP, CUDA, CuPy and JAX-GPU are unavailable. \textbf{pyopencl works} (OpenCL 2.0, int64 atomics, fp64).
\item One work-item per individual draws its sex and survival; one per offspring draws its parents.
\item Draw parents by \emph{rejection sampling among survivors}, which is exactly uniform.
\item Use a counter-based hash RNG (splitmix64 of (seed, generation, individual, stream)), so results are reproducible without per-thread state.
\item Validate statistically, not bit-for-bit: KS tests on 400 replicate populations per comparison (\code{tests/test_gpu.py}).
\item Measured throughput: CPU (numba, 8 threads) $\approx3.8\times10^7$ individual-generations/s; RX 560X with large ensembles $2.7$--$3.8\times10^8$, i.e.\ 7--10$\times$.
\item Small ensembles are limited by kernel-launch overhead (about $6\times10^7$).
\end{itemize}

\section{Step 7: ML layer (analysis only)}
Nothing from this layer feeds back into the \IBM.
\begin{description}
\item[GP denoiser (M1).] Heteroscedastic GP on $\tanh(\tfrac12\log\text{ratio})$ of \IBM\ invasion outcomes. Three settings are required (\S\ref{tr:gp}):
  \begin{itemize}
  \item a variance floor;
  \item a length-scale lower bound equal to the grid spacing;
  \item the bounded tanh response.
  \end{itemize}
  The emergent region is compared with the oracle by intersection-over-union (IoU).
\item[Neural emulator (M2).] An MLP trained on \IBM\ fixation data for $N\le600$, tested at $N=800$.
  \begin{itemize}
  \item Inputs $(\log N,s)$ do not extrapolate.
  \item The diffusion scaling variable $Ns$ does (RMSE 0.010 vs Kimura): the data obey Kimura's $Ns$ collapse.
  \end{itemize}
\item[Inverse problem (M3).] Recover $(s_f,s_m)$ from one trajectory, with an amortised MLP or an oracle least-squares fit. The difference $s_f-s_m$ is better identified than either coefficient.
\end{description}
